\documentclass[12pt, fleqn, openany]{studyGuide}
\setstretch{1.4}
\enableInlineReview
\renewcommand{\VMBookTitle}{Audit Review}
\renewcommand{\VMBookSubtitle}{Grade 7 --- Ch Graphic Review --- Changed Questions}
\renewcommand{\VMFooterURL}{https://viewmath.com/grade7}
\renewcommand{\VMFooterURLDisplay}{ViewMath.com/Grade7}
\begin{document}

\begin{center}
\begin{tcolorbox}[
	enhanced,
	colback=white,
	colframe=funTeal!60,
	boxrule=1.5pt,
	arc=10pt,
	outer arc=10pt,
	width=0.9\linewidth,
	halign=center,
	top=5mm, bottom=5mm,
]
	{\Large\bfseries\sffamily\color{funTealDark}Audit Review}\\[2mm]
	{\large\sffamily\color{funGrayDark}Grade 7 --- Ch Graphic Review --- Changed Questions}\\[2mm]
	{\normalsize\sffamily\color{funGrayDark}Verify corrections below}
\end{tcolorbox}
\end{center}

\newpage
\section*{m-9-3-q23 \hfill \normalsize\texttt{tests\_questions\_bank\_2/topics\_modified/ch09-03-experimental-probability.tex}}
\begin{practiceQuestion}{m-9-3-q23}{gmc}
\begin{questionText}
The bar graph shows the results of a simulation comparing experimental and theoretical frequencies of an event across different trial counts.

\begin{center}
\begin{tikzpicture}
  \begin{axis}[
    ybar,
    bar width=11pt,
    ymin=0,
    ymax=0.7,
    width=9.2cm,
    height=5.7cm,
    ylabel={Probability},
    xlabel={Number of Trials},
    symbolic x coords={$20$, $50$, $100$, $500$},
    xtick=data,
    ytick={0,0.2,0.4,0.6},
    enlarge x limits=0.24,
    axis line style={draw=black},
    tick style={draw=black},
    tick label style={font=\small},
    label style={font=\small},
    legend style={
      at={(0.5,1.03)},
      anchor=south,
      legend columns=2,
      draw=black,
      fill=white,
      font=\small,
      cells={anchor=west}
    },
    nodes near coords,
    every node near coord/.append style={font=\scriptsize, yshift=2pt, text=black}
  ]
  \addplot[draw=black, fill=funBlue!75] coordinates {($20$,0.55) ($50$,0.44) ($100$,0.38) ($500$,0.34)};
  \addplot[draw=black, fill=white] coordinates {($20$,0.33) ($50$,0.33) ($100$,0.33) ($500$,0.33)};
  \legend{Experimental, Theoretical}
  \end{axis}
\end{tikzpicture}
\end{center}

Which conclusion does the graph support?
\end{questionText}
\choiceA{More trials make the experimental probability less accurate.}
\choiceB{The theoretical probability changes with more trials.}
\choiceC{As the number of trials increases, the experimental probability approaches the theoretical probability.}
\choiceD{$20$ trials gives the most accurate result.}
\correctAnswer{C}
\explanation{The experimental probability drops from $0.55$ at $20$ trials to $0.34$ at $500$ trials, getting closer to the theoretical $0.33$. This demonstrates the Law of Large Numbers.}
\end{practiceQuestion}

\newpage
\section*{9-5-q22 \hfill \normalsize\texttt{tests\_questions\_bank\_2/topics/ch09-05-sample-spaces-for-compound-events.tex}}
\begin{practiceQuestion}{9-5-q22}{gmc}
\begin{questionText}
Look at the tree diagram below showing the outcomes of picking a snack and a drink.

\begin{center}
\begin{tikzpicture}[x=1cm, y=1cm, >=Latex, font=\small]
  \node at (0,4.2) {Start};
  \node at (-3,2.7) {Apple};
  \node at (0,2.7) {Granola};
  \node at (3,2.7) {Crackers};
  \node at (-4,0.9) {Water};
  \node at (-2,0.9) {Juice};
  \node at (-1,0.9) {Water};
  \node at (1,0.9) {Juice};
  \node at (2,0.9) {Water};
  \node at (4,0.9) {Juice};

  \draw[thick,->] (0,3.95) -- (-3,3.02);
  \draw[thick,->] (0,3.95) -- (0,3.02);
  \draw[thick,->] (0,3.95) -- (3,3.02);

  \draw[thick,->] (-3,2.42) -- (-4,1.18);
  \draw[thick,->] (-3,2.42) -- (-2,1.18);
  \draw[thick,->] (0,2.42) -- (-1,1.18);
  \draw[thick,->] (0,2.42) -- (1,1.18);
  \draw[thick,->] (3,2.42) -- (2,1.18);
  \draw[thick,->] (3,2.42) -- (4,1.18);
\end{tikzpicture}
\end{center}

How many outcomes are in the sample space?
\end{questionText}
\choiceA{$3$}
\choiceB{$5$}
\choiceC{$6$}
\choiceD{$9$}
\correctAnswer{C}
\explanation{The tree diagram shows $3$ snack choices each paired with $2$ drink choices. Counting the branch endpoints gives $3 \times 2 = 6$ outcomes.}
\end{practiceQuestion}

\newpage
\section*{9-5-q24 \hfill \normalsize\texttt{tests\_questions\_bank\_2/topics/ch09-05-sample-spaces-for-compound-events.tex}}
\begin{practiceQuestion}{9-5-q24}{gmc}
\begin{questionText}
A tree diagram shows outcomes for choosing a pen color (Black or Blue) and a pen type (Gel, Ballpoint, or Felt-tip).

\begin{center}
\begin{tikzpicture}[x=1cm, y=1cm, >=Latex, font=\small]
  \node at (0,4.4) {Start};
  \node at (-2.8,2.8) {Black};
  \node at (2.8,2.8) {Blue};
  \node at (-4.6,1.0) {Gel};
  \node at (-2.8,1.0) {Ballpoint};
  \node at (-1.0,1.0) {Felt-tip};
  \node at (1.0,1.0) {Gel};
  \node at (2.8,1.0) {Ballpoint};
  \node at (4.6,1.0) {Felt-tip};

  \draw[thick,->] (0,4.12) -- (-2.8,3.08);
  \draw[thick,->] (0,4.12) -- (2.8,3.08);

  \draw[thick,->] (-2.8,2.52) -- (-4.6,1.28);
  \draw[thick,->] (-2.8,2.52) -- (-2.8,1.28);
  \draw[thick,->] (-2.8,2.52) -- (-1.0,1.28);
  \draw[thick,->] (2.8,2.52) -- (1.0,1.28);
  \draw[thick,->] (2.8,2.52) -- (2.8,1.28);
  \draw[thick,->] (2.8,2.52) -- (4.6,1.28);
\end{tikzpicture}
\end{center}

Which outcome is NOT shown in this tree diagram?
\end{questionText}
\choiceA{Black Gel}
\choiceB{Blue Ballpoint}
\choiceC{Red Felt-tip}
\choiceD{Blue Gel}
\correctAnswer{C}
\explanation{The tree diagram only includes Black and Blue pen colors. Red is not one of the options, so ``Red Felt-tip'' is not in the sample space.}
\end{practiceQuestion}

\end{document}